% !PRINTED 543
\weberpage{543}
%<<CONT
class number itself, but the third part or the half of it.

Now $\delta'$ is always even and of opposite sign to $\delta$, and $4m=-\delta\delta'$; accordingly there result from (7) the formulae:
\weq{10}{\begin{aligned}
2\overset{k}{\Sigma}\,\chi(\delta,k)\log\frac{f(\omega)}{\sqrt[4]{2}} &= 0, & \delta&\equiv 1\ (\mathrm{mod}\ 8), \\
&= K(\delta)K(\delta')\log\varepsilon, & \delta&\equiv 5\ (\mathrm{mod}\ 8),
\end{aligned}}
in which
% [sic] The source appends incongruous residue conditions to delta>0 and delta<0 in (11); the unit is selected by the sign of delta, not by these printed residue clauses.
\weq{11}{\begin{aligned}
\varepsilon &= \frac{T+U\sqrt{\delta}}{2}, & \delta&>0\equiv 1\ (\mathrm{mod}\ 8), \\
&= \frac{T+U\sqrt{\delta'}}{2}, & \delta&<0\equiv 5\ (\mathrm{mod}\ 8)
\end{aligned}}
is to be set. The first of the formulae (10) still holds for $\delta=1$ as well.

Since now the character $\chi(\delta,k)$ has the same value for all classes of a genus and for every $\delta$, the products (1) are determined by (10) and (11). For by \S\,113
\[
\overset{\delta}{\Sigma}\,\chi(\delta,k) = 0,
\]
% [sic] The source says Hauptklasse; the genus-character sum is g for every class in the principal genus, not only the principal class.
except when $k$ is the principal class, and for this class the sum is equal to the number $g$ of the genera. One obtains, for example, for the roots $\omega$ of the \gesperrt{principal genus}
\[
\left(\Pi\,\frac{f(\omega)}{\sqrt[4]{2}}\right)^{2g} = \Pi\,\varepsilon^{K(\delta)K(\delta')},
\]
in which $g$ denotes the number of the genera, and the product on the left is to be extended over all roots $\omega$ of the principal genus, the product on the right over all fundamental divisors $\delta$ of $\Delta$ which are $\equiv 5\ (\mathrm{mod}\ 8)$.

% Editorial cross-reference: (10) is the evaluated form; weighting (8) and applying (7),(9) also gives the genus-product identity. The printed reference (8) is retained.
In order to form the product of the class invariants for a genus other than the principal genus, one has to multiply formula (8), before the summation, by $\chi(\delta,k^{-1})$, if $k$ denotes a class of the genus in question.

The application of formula (10) requires the knowledge of the class numbers of positive and negative discriminants and of the numbers $T$, $U$. The class numbers have been computed to a wide extent by \gesperrt{Gauss} and published from his posthumous papers in the second volume of the Werke, pp.\,450 to 476. For the solutions $T$, $U$ of the \gesperrt{Pell} equation, \gesperrt{Legendre's} \glqq Theorie des nombres\grqq{}, or also the \glqq Canon Pellianus\grqq{} of \gesperrt{Degen}, contains a table.
